
本章では、本研究にて実施した実験の方法について述べる。

\section{実験参加者}
実験参加者は合計18名（男性11名、女性7名）で、年齢は21歳から26歳（平均年齢 22.7歳、標準偏差 1.37）であった。全ての参加者は聴力および視力（矯正視力含む）に異常がないことを自己申告により確認し、実験の目的および手順、安全性のリスク（VR酔い等）について十分な説明を受けた上で、書面による同意を得て実験に参加した。
実験結果に影響を与える可能性のある個人の属性として、音楽経験およびVRデバイスの使用経験について事前アンケートを実施した。音楽経験に関しては、「楽器演奏経験が3年以上ある」と回答した経験者が9名、「半年未満（または経験なし）」と回答した未経験者が9名となり、両群が同数となるよう構成された。また、相対音感トレーニングの専門的な受講経験については、16名が「経験なし」と回答し、過去に受講経験がある者は2名のみであった。 　VRデバイスの使用経験に関しては、「日常的に使用している」または「慣れている」「頻繁ではないが慣れている」と回答した者が計13名（約72%）を占め、「数回使用したことがある」が4名、「初めて」が1名であった。このことから、大半の被験者はVR空間での操作に対して一定の順応性を有していたと考えられる。


\section{実験手順}
本実験の所要時間は被験者1名につき約30分である。実験の開始に先立ち、被験者に対して実験概要の説明を行い、現在の体調や聴覚に関する1〜2分程度の事前アンケートを実施する。その後、被験者はHMDおよびコントローラーを装着し、実験タスクへと移行する。
本実験では、前方範囲のみを使用する180度環境と、周囲全体を使用する360度環境の2つの条件において聴取実験を行う。この際、実験を行う順序によって生じる学習効果や疲労などの順序効果を排除するため、180度環境と360度環境のどちらを先に行うかは被験者ごとにランダムに決定し、カウンターバランスをとる設計とした。また、2つの条件の間には適度な休憩を設け、被験者の状態をリセットした上で次の条件へと進むものとする。
各条件における実験は、まず約5分間の練習を行い、被験者を各環境特有の空間音響や操作方法に十分に習熟させる。続いて、全24回の音階テストを行う。テストフェーズの所要時間は約5分である。両条件の終了後、1〜2分程度の事後アンケートおよびNASA-TLXを用いたシステム負荷の評価を実施し、実験を終了する。

\section{評価手法}
\subsection{定性評価}



\subsection{3つの定量評価}


